\documentclass[12pt]{article}
\usepackage[utf8]{inputenc}
\usepackage[T1]{fontenc}
\usepackage[margin=1in]{geometry}
\usepackage{setspace}
\usepackage{xcolor}
\usepackage[colorlinks=true, urlcolor=blue!60!black, linkcolor=blue!60!black]{hyperref}
\setlength{\parindent}{0em}
\setlength{\parskip}{0.6em}
\renewcommand{\baselinestretch}{1.0}
\begin{document}
\begin{center}
{\LARGE \textit{Tennessee v. Lane}, 541 U.S. 509 (2004)\par}
\vspace{0.5em}
{\large Case Brief --- Two-Page Edition\par}
\end{center}

\section{Memory Jogger}
Title~II of the ADA validly abrogates state sovereign immunity as applied to the fundamental right of access to the courts (\href{https://www.law.cornell.edu/supremecourt/text/541/509}{541 U.S. 509}).

\section{Detailed Case Facts}
George Lane, a paraplegic summoned to a second-floor Tennessee courtroom with no elevator, crawled up two flights of stairs for one appearance; when he refused to crawl or be carried at the next, he was jailed for failure to appear. Beverly Jones, a court reporter who uses a wheelchair, lost work because of inaccessible courthouses. They sued Tennessee for damages under Title~II, which expressly abrogates state immunity.

\section{Procedural History}
The district court denied Tennessee's immunity motion; after \textit{Garrett}, the Sixth Circuit affirmed on a due process theory; the Supreme Court affirmed.

\section{Judicial Votes}
5--4. Stevens wrote for O'Connor, Souter, Ginsburg, and Breyer; Souter and Ginsburg concurred; Rehnquist (with Kennedy and Thomas), Scalia, and Thomas dissented.

\section{Holding}
Under \textit{Boerne}'s congruence-and-proportionality test, assessed as applied to the category of rights at stake, Title~II validly abrogates immunity in cases implicating access to the courts. Other settings were left open.

\section{Analysis of Opinions}
Stevens distinguished \textit{Garrett}: Title~II enforces fundamental rights subject to heightened scrutiny, making a pattern of violations easier to show, and the record documented inaccessible public buildings and courthouses. Its remedy, reasonable modifications without undue burden, is proportional. Ginsburg described equal citizenship as requiring accommodation. Rehnquist found little evidence of state violations and accused the majority of rigging the test with a courthouse-only statute; Scalia rejected the test as flabby and confined Section~5 to remedying actual violations.

\section{Examples: Future Applications}
\textbf{Same-Side:} a deaf parent who is denied an interpreter at a termination hearing may sue. \textbf{Opposite-Side:} inaccessible seating at a state hockey rink implicates no fundamental right, so \textit{Garrett} controls. \textbf{Fence-Sitter:} a university's refusal of extra exam time.

\section{Critique}
Textualists echo Scalia that ``enforce'' does not permit prophylaxis and that the test invites policymaking. Disability-rights scholars praise the result but say the piecemeal holding left plaintiffs litigating category by category.

\section{Key Quotations}
\begin{quote}
``For these reasons, we conclude that Title II, as it applies to the class of cases implicating the fundamental right of access to the courts, constitutes a valid exercise of Congress' \S~5 authority to enforce the guarantees of the Fourteenth Amendment.''
\end{quote}
\noindent --- Justice Stevens, Opinion of the Court, 541 U.S. at 533--534.

\end{document}
